\section{置换检验依赖可交换性}
\label{sec:11-Mathematical-Statistics/06-Resampling-and-Model-Selection/02}

\subsection{两条田埂、两组高矮、一个数字}

你站在稻田边上。左边田施了新肥，右边照旧，各取了 8 株稻子量高度（cm）：

\[
\begin{aligned}
\text{新肥：} &\quad 92,\; 88,\; 95,\; 90,\; 97,\; 86,\; 93,\; 91\\
\text{旧肥：} &\quad 85,\; 82,\; 88,\; 79,\; 84,\; 86,\; 80,\; 83
\end{aligned}
\]

新肥组均值 91.5，旧肥组均值 83.4。差了 8.1 cm。

第一反应很自然：做 \(t\) 检验。两组独立，算 pooled variance，带入 \(t\) 公式，翻分布表，看 \(p\) 值。标准流程，三行代码的事。

但这里你要做一个小小的暂停。\(t\) 检验依赖一条假设——数据从正态分布里抽出来，至少在大样本下近似成立。你手上只有 16 株稻子，两组各 8 株。扫一眼数据：新肥组里 97 和 86 跨了 11 cm，旧肥组从 79 到 88 跨了 9 cm。8 个点没法给你底气说``这大概正态''。

你碰到的不是数学错误——是工具链里的一个空白：有一类假设检验工具，名字叫``参数检验''，它靠你事先为数据指定一个分布族（通常就是正态），然后在这个族内部去推断参数。当数据太少先验知识又薄时，这个前提站不住。

\subsection{换个角度想：标签是从帽子里抽出来的}

你真正想问的是：这 8.1 cm 是因为新肥真的有效，还是因为施肥的时候恰好有几株本身就高的稻子分到了新肥组？

如果肥料毫无作用，每株稻子在新肥和旧肥下的高度应该一模一样。那么你看到的两个组之间的差异，纯属 16 株稻子被随机贴上了``新肥''和``旧肥''标签的结果——标签是随机的，高度是固定的。

这样一来，你就可以模拟：如果随机打乱标签，两组均值差会是什么样？

先算一遍真实的数据：16 个数摞在一起，总共 \(\binom{16}{8} = 12870\) 种标签分配方式，保持两组各 8 个。对每一种分配，算``贴上'新肥'标签的那 8 株的平均''减去剩下 8 株的平均''。得到 12870 个差值。画成直方图——这就是如果肥料无效、纯靠运气，差值能有多大。

真实的 8.1 cm 在这个直方图的什么位置？如果它落在分布的尾巴上——比如只有 3\% 的随机排列能产出 8.1 以上的差值——那你就可以说：在肥料无效的假设下，看到 8.1 是相当不寻常的。要么你的运气极其离谱，要么肥料确实有效。

这就是置换检验的核心。它不从正态分布抽 \(p\) 值——它从随机化本身抽。

选统计量

\[
T(Y, Z) = \bar Y_{\text{treated}} - \bar Y_{\text{control}},
\]

双侧 \(p\) 值为

\[
p = \frac{
\#\{\pi : |T^\pi| \ge |T_{\mathrm{obs}}|\}
}{
\#\{\text{所有符合条件的排列}\}
}.
\]

\subsection{排列太多怎么办：随机抽一个子集}

\(\binom{16}{8} = 12870\) 还能枚举。换到两组各 50 人，\(\binom{100}{50}\) 是天文数字——枚举不可行。

解决方法直接到朴素：不枚举全部，随机抽 \(B\) 个排列，用这 \(B\) 个的经验分布逼近真分布。Monte Carlo 修正公式：

\[
\hat p = \frac{
1 + \#\{b : |T^{\pi_b}| \ge |T_{\mathrm{obs}}|\}
}{
B + 1
}.
\]

分子分母各加一，把真实观测也纳入零分布。如果原始观测的差异在 \(B\) 次随机排列中都排第一，不加修正会输出 \(p = 0\)——但有限模拟不可能支撑精确零概率。加一之后 \(p\) 至少是 \(1/(B+1)\)，与``我们只试了 \(B\) 种排列''的精度匹配。

\subsection{最关键的不是算 \(p\)——是凭什么可以打乱}

读到这里，置换检验给人的第一印象很可能是``万能''——不需要正态，不需要等方差，甚至不需要大样本。只要你有计算机，把它往排列上一扔，就算出了一个精确的 \(p\) 值。许多教材把置换检验归入``非参数检验''一类，这一类名字本身就在暗示``把参数假设丢掉了，更自由了''。

这里有一个需要非常小心对待的细节。丢掉了正态假设，不等于丢掉了所有假设。置换检验用另一个假设换了参数假设的位置——可交换性。

置换合法，是因为在原假设下，排列不改变数据的联合分布。形式化地说：如果 \(H_0\) 为真，且随机变量序列 \(Y_1, \ldots, Y_n\) 在标签的某种排列群下分布不变，那么打乱标签给出的参考分布就是精确的。

可交换性必须匹配实验设计，一对一套牢：

\begin{itemize}
\tightlist
\item 配对实验中，只能在配对内部交换两个标签或做符号翻转——不能把张三的测量值贴到李四头上；
\item 分层随机化中，置换只能局限在每个层内部——跨层打乱破坏了层的结构；
\item cluster-randomized trial 中，以整簇为单位整体置换——把簇内成员拆散是对随机化机制的扭曲；
\item 时间序列中，时序本身携带着依赖信息——任意打乱就是在假设独立性，而时间序列从一开始就不是独立的；
\item 不等概率设计中，需要按实际随机化方案的分布去生成排列，均匀随机地打乱是错的。
\end{itemize}

``非参数''绝不等于``无假设''。参数检验假设分布形状（正态、指数等等）；置换检验假设一种对称性或者随机化结构。它们是不同种类的假设，一条换成了另一条，不是一条消失了。

\subsection{Sharp null 与平均效应是两个原假设}

这一步的坑更深。置换检验在 sharp null 下才是精确的。

Sharp null 说：每一个实验单位的处理效应都是零。张三无论分到处理组还是对照组，他的响应值完全一样。在这个假设下，每个单位的潜在结果只有一种固定数值——你可以为任何标签排列填出相应的``观测值''，置换因而合法。

``平均处理效应为零''是另一个完全不同的原假设。它允许张三的处理效应为 \(+10\)、李四为 \(-10\)，两人一平均刚好为零。在这个假设下，张三的响应值在两组中不一样——他的数据确实受处理影响了，只是平均起来看不见。此时打乱标签并不能保持联合分布，置换产生的 \(p\) 值不再具有精确校准。

这两类原假设下的置换分布不同。你用全排列分布去检验``平均效应为零''，检验尺可能偏或歪。可以用学生化统计量（如 studentized mean difference）和渐近理论来改善，但必须清楚自己正在检验什么命题。

这里只是借用潜在结果框架来区分两种零假设。要系统理解平均处理效应如何定义、随机化为何能识别它、观测研究需要什么额外条件，可继续读 \hyperref[ch:120]{``因果推断：预测世界不等于改变世界''}。

两样本的标签置换精确检验的是``两组分布完全一致''。如果目标只是检验均值相等而方差不一致——用未学生化的均值差做置换，方差不同的两组标签排列就已经改变了数据的结构，检验可能失准。

\subsection{Bootstrap 和 permutation 各做各的事}

两个方法都涉及``重新摆弄数据''，容易搞混，但推理逻辑完全不同。

Bootstrap 从数据的经验分布中有放回地再抽样。它模拟的是：如果再从同一个总体里抽一批新数据，估计量会波动多大。Bootstrap 样本会重复某些观测、跳过另一些观测，以此来模仿抽样变异性。它的本职是标准误和置信区间。

Permutation 在原假设的对称性下重新分配标签。它模拟的是：如果随机化机制跑出了另一个结果，检验统计量会是多少。它只重新排列已有的标签，不改动观测值本身。它的本职是检验。

Bootstrap 给你一个估计量的波动尺度；permutation 给你一个零假设下的参考分布。把直接置换的那行抽样代码换成 bootstrap，得到的结果不能同时承担两者的解释。

\subsection{选了统计量就绑定了功效}

置换检验不对第一类错误有精确控制——它确实做得很好。但这个保证不送功效。功效来自统计量跟备择假设的匹配。

均值差只擅长检测位置平移。如果处理的效应不在中心而在尾部——新肥没有让每株稻子都高一点，而是只让少数稻子猛窜、大部分不变——用均值差做置换检验的功效会很低。因为多数排列中，均值差对尾部的少数极端值不敏感，真实的尾部效应被埋在大量不变的数据点里。

可以针对不同的备择选不同的统计量：秩和统计量对位置平移之外的单调效应有更好的稳健性；energy distance 或者 maximum mean discrepancy 可以捕捉更一般的分布差异。但统计量是在看到数据之前就应该定下来的。如果你看了数据，试了均值差、中位数差、秩和，然后挑了 \(p\) 值最小的那个来报告——置换过程必须重现整个选择步骤，或者用独立数据做选择。只在最终选中的统计量上做一次置换，忽略选择的多重性，\(p\) 值会虚低。这不是置换的错，是多重比较的错，任何检验框架里都一样。

置换检验最有力的地方不在``甩掉假设''——而在``用随机化机制代替分布形状假设''。这个置换有条件：设计匹配的可交换性、sharp null 的范围、统计量和备择假设的契合。承认这些条件，才能把置换从万能错觉里摘出来，放进恰好适合它的位置。
